% Source: sources/returns/S6_04_MAPS_MANIFOLDS/sections/manifolds.tex:5448–5782
\paragraph{Theorem D.2.27 (Proved: the three $I_1^*$ profiles give no natural sphere).}
None of the natural section-line inputs on
\[
 (I_1^*,I_3,I_1,I_1),\qquad (I_1^*,II,I_3),\qquad
 (I_1^*,IV,I_1)
\]
produces a homotopy six-sphere, for any permitted original,
logarithmic or Mumford filling, multiplicity or affine twist in
the sector of D.2.22. The obstruction is integral. In the only
remaining minimal star caps, the full normal-kernel indices are
\[
\begin{array}{c|c|c|c}
 m_*&\text{component action}&H_2(U_*;\mathbb Z)
       &\kappa_*\ \text{(normal projection index)}\\\hline
 2&\text{one component}&\mathbb Z^2&1\\
 4&\text{two components interchanged}&
                   \mathbb Z^2\oplus\mathbb Z/2&4.
\end{array}                                                   \tag{D.ST1}
\]
If the complete output is simply connected and its other caps
have the minimal rational ranks, the order of $H_2$ has exactly
one factor of $2$ in the first row, and exactly two factors of $2$
in the second row. In particular $H_2\ne0$.

\begin{proof}
\emph{1. Exhausting the possible star caps.}
The sixteen natural initial rows on these three profiles are
listed in the certificate
\path{records/manifolds_i1star_obstruction.json}.
No original star cap has both zero component multidegree and one
component orbit, so D.2.22 excludes every such original cap from
a sphere. A semistable root cover has order $m_*=2k$, and
$A_*^{2k}-1=kN$, where $N$ has rank one, square zero and primitive
image. Thus its minimal rank-one fibre has exactly $k$ components.
This is the integral version of the $I_n^*$ semistable reduction
used in D.2.24; see also \cite[Sections 7.6--7.7]{manSS10}.

On the dual-complex circle the residual action reverses
orientation. Its action on the $k$ components therefore has the
form $i\mapsto r-i$ and has at least $\lceil k/2\rceil$ orbits.
The transverse two-torus has trivial linear action after the
vanishing circle is collapsed. The component classes contribute
one class per orbit, the transverse area contributes one, and
the dual-circle cross terms are anti-invariant. Consequently
\[
 b_2(U_*)=\rho+1,\qquad
 \dim H_2(T^4;\mathbb Q)_{A_*}=2.
\]
The saturated rank budget of D.2.22 forces $\rho=1$.
Hence $k=1$, or $k=2$ with the two components interchanged.
Further toric subdivisions add a component defect and cannot
restore this equality. This proves that (D.ST1) exhausts the
remaining cap types without an order bound imposed externally.

There is a second necessary integral condition. On the quotient
of the fibre lattice by the primitive vanishing circle, an
actual affine power vector is fixed by the induced involution.
In eight of the sixteen rows the image of this fixed lattice
under $\gamma^*$ is $2\mathbb Z$. These are precisely the ordered
section pairs
\[
 (a,b)=(4,1),\ (4,3),\ (12,1).
\]
Their star numerator is even, while their star multiplicity is
even. The Seifert presentation D.2.23 therefore precludes simple
connectivity. This uses the fixed lattice \emph{after} collapsing
the vanishing circle, not the smaller fixed lattice before that
collapse.

For the other eight rows the pairs are
\[
 (a,b)=(1,4),\ (1,12),\ (3,4),                             \tag{D.ST2}
\]
and there is an integral basis $(\alpha,\beta,\lambda,\ell)$
with $\gamma^*(\ell)=1$ and
\[
\begin{gathered}
 A\alpha=-\alpha,\quad A\beta=-\beta-\lambda,\quad
 A\lambda=\lambda-\alpha,\quad A\ell=\ell,\\
 \delta=-\alpha+2\lambda,\qquad
 A^2-1=\alpha\otimes\beta^*.
\end{gathered}                                               \tag{D.ST3}
\]
Here $\alpha$ is primitive. The exact integral change of basis
for every marked row is in the certificate. For example, in a
standard marking put
\[
\begin{aligned}
 \alpha&=2e_1+a\delta,&
 \lambda&=e_1+\tfrac{a+1}{2}\delta,\\
 \beta&=e_2+\tfrac{3a-1}{2}\lambda,&
 \ell&=-\tfrac{3b}{4}e_1+\tfrac b2e_2+\gamma.
\end{aligned}
\]
The exceptional marking on $50/II$ is also checked integrally.
In this basis the degree-two fibre coinvariants are freely
generated by
\[
 x=[\beta\wedge\lambda],\qquad y=[\beta\wedge\ell],
\]
with relations
\[
 [\alpha\wedge\beta]=-2x,\quad
 [\lambda\wedge\ell]=-2y,\quad
 [\alpha\wedge\lambda]=[\alpha\wedge\ell]=0.                \tag{D.ST4}
\]
In particular the boundary group $H_2(\partial U_*;\mathbb Z)$
is free of rank four; there is no boundary torsion that can be
discarded. The exterior fibre functional is still
$\omega=a e_1^*e_2^*+b\gamma^*\delta^*$, and
\[
 \omega(x)=-a,\qquad \omega(y)=b/4.                         \tag{D.ST5}
\]

\emph{2. The normalization and its integral parity.}
We compute the retraction of $U_*$ to its central fibre by its
normalization and conductor. Since a simply connected output
has $\gcd(m_*,n_*)=1$, its star numerator is odd. All quotient
maps on the transverse elliptic curves below are therefore
unramified, with the displayed degrees.

For $m_*=2$ the normalization is an oriented $S^2$ bundle $N$
over a two-torus. Its conductor is a connected bisection $C$,
which maps with degree two to the double curve $D$. The relation
$A\lambda=\lambda-\alpha$ in (D.ST3) determines the parity of
this bundle. In a transverse circle coordinate $u$ its
pole-interchanging clutch can be represented by
\[
 z\longmapsto e^{-2\pi i u}/z.
\]
The corresponding $\mathrm{PGL}_2$ loop has representative
$\left(\begin{smallmatrix}0&e^{-2\pi i u}\\1&0\end{smallmatrix}\right)$.
Its determinant has odd winding, so its lift to the double cover
of $\mathrm{SO}(3)$ changes sign. Hence $w_2(N)$ evaluates to one
on the base. The remaining affine translations and unit
multipliers change the flat part of the clutch, not this winding;
the odd numerator changes the primitive base marking by an odd
index and preserves the parity. This connects the parity
calculation to every actual affine cap, not only a selected
linear model.

Let $f$ be the sphere fibre and choose a section class $\sigma$
with $\sigma^2=1$. Such a choice follows by adding multiples of
$f$ to any odd-square section. The bisection has zero
self-intersection, so
\[
 [C]=2\sigma-f.
\]
The normalization pushout gives
\[
 H_2(U_*)=
 \langle\sigma,f,D\mid 2\sigma-f-2D=0\rangle
          \cong\mathbb Z^2.                               \tag{D.ST6}
\]
There is no additional $H_1$-kernel term: the conductor map to
$D$ is injective on $H_1$. The regular torus
$\alpha\wedge\beta$ maps to $f$, since it sweeps the single
semistable component with degree one. The regular
$\lambda\wedge\ell$ torus lies in the smooth part of the
normalization, covers its base twice and is disjoint from $C$.
Its class is consequently $2\sigma-f$, hence $2D$ in the
pushout. By (D.ST4),
\[
 x\longmapsto-\sigma+D,\qquad y\longmapsto-D.
\]
These are an integral basis. Thus the fibre-coinvariant map has
cokernel zero in this case.

For $m_*=4$ with the two components interchanged, choose one
component upstairs. Its stabilizer is generated by the square
of the deck generator. That square fixes
$\alpha,\lambda,\ell$. The open $\mathbb C^*$ bundle of its
normalization therefore has an abelian period lattice. More
explicitly it has generators $\alpha,\lambda,\ell,c^2$ and one
relation
\[
 2c^2=t\alpha+r\lambda+n_*\ell.
\]
Because $n_*$ is odd, $\alpha$ is primitive in this lattice.
The associated complex line over the quotient two-torus has
zero Chern class. Its two pole sections are thus homologous;
write their common class as $\sigma$. They map with degree two
to two different double curves $D_0,D_1$. Hence
\[
\begin{split}
 H_2(U_*)&=
 \langle f,\sigma,D_0,D_1\mid
                  \sigma-2D_0=\sigma-2D_1=0\rangle\\
 &\cong\mathbb Z f\oplus\mathbb Z D_0
                       \oplus\mathbb Z/2\,(D_1-D_0).
\end{split}                                                 \tag{D.ST7}
\]
Again the conductor map on $H_1$ is injective into the sum of
its double-curve targets, so this is the complete group.

The regular $\alpha\wedge\beta$ torus sweeps both semistable
components, each with degree one onto the chosen component
fibre downstairs, and therefore maps to $2f$. The regular
$\lambda\wedge\ell$ torus covers the normalization base twice.
It avoids both pole sections, so its class is $2\sigma=4D_0$.
Thus, modulo possible multiples of the displayed order-two
class,
\[
 x\longmapsto-f,\qquad y\longmapsto-2D_0.
\]
The free part of this map has determinant two. Including the
order-two summand, its full cokernel has order four, independent
of the two possible torsion corrections. This is precisely the
integral effect missed by the rational vertical coinvariant
identity $[\delta]=4[\beta]$.

For completeness, the conductor has no third homology, and
its degree-two map into the direct sum of the double-curve
groups is injective: each conductor component maps with degree
two to its corresponding double curve. Thus the normalization
sequence identifies $H_3(U_*)$ with $H_3(N)=\mathbb Z^2$.
It also gives $H_4(U_*)=H_4(N)=\mathbb Z$, generated by the
reduced central component. A positive multiple of this divisor is principal,
being the divisor of the base parameter. Its intersection with
every two-cycle is therefore zero. Under Poincare--Lefschetz
duality this is the entire map
$H_2(U_*)\to H_2(U_*,\partial U_*)=H^4(U_*)=\mathbb Z$.
It follows that
\[
 H_2(\partial U_*)\longrightarrow H_2(U_*)
 \quad\hbox{is onto integrally}.                            \tag{D.ST8}
\]

Let $C_*=H_2(T^4)_{A_*}$ and let $K_*$ be the full kernel in
(D.ST8). The boundary Wang sequence gives
\[
 0\longrightarrow C_*\longrightarrow H_2(\partial U_*)
        \longrightarrow\operatorname{Fix}(A_*)
        \longrightarrow0.
\]
The fibre map $C_*\to H_2(U_*)$ is injective, and the preceding
calculations therefore give
\[
 [\operatorname{Fix}(A_*):\operatorname{proj}K_*]
 =\bigl|H_2(U_*)/\operatorname{im}C_*\bigr|
 =\begin{cases}1,&m_*=2,\\4,&m_*=4\text{ with swap}.
   \end{cases}                                             \tag{D.ST9}
\]
This proves (D.ST1), retaining the normalization torsion and
both full normal directions.

\emph{3. The full normal determinant and the global obstruction.}
Here is a useful integral determinant calculation. Suppose all
caps have the minimal rational ranks and their boundary maps
in degree two are onto. Put
\[
 \kappa_i=[\operatorname{Fix}(A_i):
                         \operatorname{proj}K_i],
\]
where $K_i$ is the full integral boundary kernel. The primitive
graph calculation for all $172$ rows gives the exterior group
as $\mathbb Zs$ plus the direct sum of the fixed sweep lattices,
with the sole relation $\sum_i t_i(\delta)+be\,s=0$.
Over $\mathbb Q$, each $K_i$ is the graph of a linear fibre
functional on its fixed lattice. Its value on $\delta$ is fixed
without choosing any of its other coefficients: the
$\delta$ normal torus bounds on the root cover, and transfer
with the actual affine power vector gives the exact relation
\[
 m_i t_i(\delta)+bn_i s=0.                                 \tag{D.ST10}
\]
The sign is the same oriented convention as (D.HY3).

Choose integral bases of the projected kernels. Their
covolumes are $\kappa_i$. Eliminate all independent sweep
coordinates over $\mathbb Q$ and then the single diagonal
$\delta$ relation. The determinant of the complete square
Mayer--Vietoris matrix is
\[
 \left(\prod_i\kappa_i\right)
                   b\left(e-\sum_i n_i/m_i\right).         \tag{D.ST11}
\]
All longitudinal and mixed fibre coefficients cancel: only
the values (D.ST10) on the diagonal vector enter the last
Schur complement. Multiplying the rational calculation by the
integral covolumes is an equality of integer determinants,
not a saturation assumption. The certificate also verifies
this identity symbolically with arbitrary two-by-two normal
projection matrices and arbitrary longitudinal coefficients.

For an unchanged minimal cap or an index-one multiplicative
cap, $\kappa_i=1$. For a finite good-reduction cap D.2.26 gives
$\kappa_i=m_i/q_i$. For a regular logarithm its projected
kernel is $m_i\Lambda+\mathbb Z v_i$, of index $m_i^3$;
this is the full lattice because the affine action is free.
These statements include every cap that could remain in a
sphere after the D.2.22 rank obstruction.

If $X$ is simply connected, the Seifert determinant is $\pm1$,
so
\[
 \left|e-\sum_i n_i/m_i\right|=(\prod_i m_i)^{-1}.
\]
All degree-one groups in this Mayer--Vietoris map are free.
For the star boundary this follows from (D.ST3); for the
remaining boundaries the local matrices have only unit
nonzero Smith factors, as checked in every retained row.
The star cap has the abelian presentation obtained from the
collapsed lattice and the relation $m_*c=v_*$; its relation
is primitive because $\gcd(m_*,n_*)=1$. The same primitive
relation proves freeness for the finite and regular log caps.
An unchanged minimal additive cap has the product/mapping-torus
model of D.2.26 and free first homology. The exterior group is
free as well. The degree-one Mayer--Vietoris source and target
have equal rank, while their cokernel is $H_1(X)=0$.
Consequently that map is an isomorphism and has no additional
degree-one bridge kernel.
Thus (D.ST11) gives the exact order
\[
 |H_2(X)|=\frac{b}{q_*Q_o}\,R^2,\qquad
 q_*:=m_*/\kappa_*=
 \begin{cases}2,&m_*=2,\\1,&m_*=4\text{ with swap},\end{cases} \tag{D.ST12}
\]
where $Q_o$ is the product of the indices $q_i$ of the other
finite caps and $R$ is the product of regular logarithmic
multiplicities.

Simple connectivity permits at most two multiple fibres and
forces their multiplicities to be coprime. Since $m_*$ is
even, every other multiplicity is odd. The only possible
additional finite index is $3$ from the $IV$ fibre of $50/IV$;
an order-six $II$ log cannot coexist with the star log in a
simply connected output. In every row of (D.ST2),
$b/Q_o$ has precisely two factors of $2$, and $R$ is odd.
Equation (D.ST12) now has two-primary order $2$ for $m_*=2$
and $4$ for $m_*=4$. Neither can be a homology sphere.
All excluded cap choices already have a positive rational
second Betti number or nontrivial fundamental group. This
completes the exhaustion.
\end{proof}
